\chapter{Le central risk book}\label{ch:s2:the-central-risk-book}

Dix desks qui couvrent chacun leur propre risque paient dix fourchettes pour des transactions qui se compenseraient peu. Un \omterm{def:s2:the-central-risk-book:crb}{central risk book} recueille d'abord le risque et croise ce qui
se compense en interne, et ne couvre que le solde, lorsque cela est peu coûteux. Sur la banque synthétique de ce chapitre, cinq desks reçoivent des flux clients indépendants sur
100 actions paieraient 335\,000~\$ par jour pour couvrir chaque flux d'un coup. Mutualisés, les flux nets coûtent 270\,000~\$, soit 19.5~\% de moins. Conservés et écoulés à 15~\% par jour, avec
le risque de marché couvert dans des contrats à terme sur indices, le dispositif coûte 45\,000~\$ par jour. Le book conserve alors un risque résiduel dont la valeur en risque à un jour et à 99~\% s'élève à 859\,000~\$.
La construction s'appuie sur \texttt{firm.crbook}.

\section{Mutualisation et compensation}

\begin{definition}[Central risk book]\label{def:s2:the-central-risk-book:crb}
Un \emph{central risk book}\index{central risk book} est un portefeuille de négociation qui prend en charge le risque que plusieurs desks acquièrent auprès de clients, le compense et
en gère le solde~: il croise en interne les positions qui se compensent, couvre les facteurs communs avec des instruments liquides et écoule le reste au fil du temps dans les
limites de risque.
\end{definition}

\begin{definition}[Mutualisation des risques]\label{def:s2:the-central-risk-book:pooling}
La \emph{mutualisation des risques}\index{mutualisation des risques} (risk pooling) désigne le regroupement de positions issues de plusieurs sources avant couverture, de sorte que les expositions opposées s'annulent et que seul le
solde soit négocié.
\end{definition}

\texttt{firm.crbook} fournit à cinq desks des flux clients sur 100 actions pendant 500 jours (voir le \cref{lst:s2:the-central-risk-book:book}). Le flux de chaque desk sur chaque valeur
a un écart type de 0.5~million de dollars par jour, et 10~\% de la variance est commune à tous les desks, comme lorsque les clients suivent le même mouvement. Les rendements des actions suivent un
modèle factoriel à trois facteurs~: un facteur de marché de volatilité quotidienne de 1~\%, deux facteurs de style et 1.5~\% de volatilité idiosyncratique. Les coûts de négociation reprennent le modèle du livre~7
(chapitre~27)~: une demi-fourchette de 3 points de base et un impact en racine carrée de coefficient 0.7, sur des volumes quotidiens compris entre 20~\$ et 200~millions de dollars.
Almgren, Thum, Hauptmann et Li, à partir de près de 700\,000 ordres, ont trouvé que l'impact temporaire est plus proche d'une puissance 3/5 du taux de négociation que d'une racine carrée~;
l'exposant compte moins ici que le fait que l'impact croît plus vite que la taille.

Couverts desk par desk, chaque flux étant exécuté isolément comme si le desk était seul sur le marché, les couvertures coûtent 335\,000~\$ par jour. Mutualisé au préalable, 47.0~\% du
flux brut se croise en interne et le solde coûte 270\,000~\$, soit une économie de 19.5~\%. L'économie est plus faible que la part croisée, car l'impact par
dollar est plus faible sur les petites transactions~: ce qui se compense en interne est disproportionnellement la part bon marché. Butz et Oomen ont décrit le choix que fait chaque dealer entre
internaliser le risque d'un client, le détenir en anticipant un flux opposé, et l'externaliser en le couvrant immédiatement~; un \omterm{def:s2:the-central-risk-book:crb}{central risk book} fait
ce choix pour l'ensemble de la banque.

\section{Croisement interne}

L'économie dépend du degré d'indépendance des flux des desks (voir la \cref{fig:s2:the-central-risk-book:desks}). Avec des flux totalement indépendants, mutualiser dix desks
permet d'économiser 48.0~\% du coût. Avec 10~\% du flux commun à tous les desks, l'économie culmine à 19.7~\% avec six desks puis tombe à 17.1~\% avec dix~: la
part commune ne se compense pas, et avec davantage de desks elle représente une part plus grande du solde. Les desks d'une banque voient les mêmes humeurs de leurs clients, et l'intérêt de la mutualisation doit
être apprécié à partir de leur corrélation effective.

\begin{omfigure}
\begin{tikzpicture}
\begin{axis}[width=11cm,height=5.2cm,xlabel={\small nombre de desks mutualisés},ylabel={\small coût de couverture économisé (\%)},
  tick label style={font=\footnotesize},xmin=1,xmax=10,ymin=0,ymax=50,grid=major,grid style={gray!25},table/col sep=comma,
  legend style={font=\scriptsize,at={(0.97,0.03)},anchor=south east}]
\addplot[omThm,thick,mark=*,mark size=1.2pt] table[x=desks,y=independent]{figdata/strategies-2/25-the-central-risk-book/desks.csv};
\addplot[omDef,thick,mark=square*,mark size=1.2pt] table[x=desks,y=common]{figdata/strategies-2/25-the-central-risk-book/desks.csv};
\legend{flux indépendants,10~\% de la variance commune}
\end{axis}
\end{tikzpicture}

\omcaption{Coût de couverture économisé en mutualisant les flux des desks avant de les couvrir d'un coup, en fonction du nombre de desks, pour des flux indépendants et pour des flux
avec une composante commune. Données~: \texttt{s2\_crbook.by\_desks}.}\label{fig:s2:the-central-risk-book:desks}
\end{omfigure}

\section{Couverture optimisée}

\begin{definition}[Couverture du risque résiduel]\label{def:s2:the-central-risk-book:residual}
Une \emph{couverture du risque résiduel}\index{couverture du risque résiduel} (residual risk hedge) couvre d'un coup, avec des instruments liquides (contrats à terme sur indices,
paniers de facteurs), l'exposition d'un book mutualisé aux facteurs communs, et laisse les positions restantes, pour l'essentiel idiosyncratiques, être écoulées progressivement ou compensées par des flux ultérieurs.
\end{definition}

Le central book n'a pas besoin de couvrir le solde d'un coup. Chaque jour, il intègre le flux net des clients à son inventaire, vend une fraction fixe de l'inventaire
en actions et détient un contrat à terme sur indice pour couvrir le reliquat pondéré par les bêtas. Garleanu et Pedersen ont montré que, lorsque la négociation est coûteuse, la politique
optimale consiste à se rapprocher partiellement de sa cible à chaque période~; ici la cible est un book neutre, et la fraction fixe le compromis. Les flux de demain compensent
une partie de ce qui est conservé, et ce qui est écoulé l'est par petites tranches~:

\begin{center}
\small
\begin{tabular}{@{}p{4.2cm}>{\centering\arraybackslash}p{3.6cm}>{\centering\arraybackslash}p{3.6cm}>{\centering\arraybackslash}p{3.6cm}@{}}
\toprule
part écoulée par jour & coût (milliers de \$ par jour) & écart type du P\&L (milliers de \$) & écart type sans contrats à terme\\
\midrule
100~\% (mutualisé, d'un coup) & 269.7 & 0 & 0\\
50~\% & 123.1 & 128 & 157\\
30~\% & 78.4 & 222 & 268\\
20~\% & 56.5 & 304 & 367\\
15~\% & 45.3 & 369 & 448\\
10~\% & 33.4 & 472 & 579\\
5~\% & 20.3 & 684 & 842\\
\bottomrule
\end{tabular}
\end{center}

La surcouche par contrats à terme coûte environ 500~\$ par jour et supprime environ un cinquième du risque~: les positions nettes sont réparties entre de nombreuses valeurs, de sorte que l'essentiel de ce qui est conservé
est idiosyncratique. Le point de la frontière (voir la \cref{fig:s2:the-central-risk-book:frontier}) qui convient dépend du prix du risque. Avec une limite de valeur en risque à un jour et à 99~\%
de 1~million de dollars, le taux le moins coûteux respectant la limite est de 15~\%~: 45\,300~\$ par jour et une valeur en risque de 859\,000~\$. Avec une limite de
500\,000~\$, il est de 35~\%~: 89\,300~\$ par jour. Par rapport à une couverture desk par desk, la première solution permet d'économiser 86.5~\% du coût.

\begin{omfigure}
\begin{tikzpicture}
\begin{axis}[width=11cm,height=5.4cm,xlabel={\small écart type du P\&L quotidien conservé (milliers de \$)},
  ylabel={\small coût de couverture (milliers de \$ par jour)},tick label style={font=\footnotesize},xmin=0,xmax=900,ymin=0,ymax=290,grid=major,
  grid style={gray!25},table/col sep=comma,legend style={font=\scriptsize,at={(0.97,0.97)},anchor=north east}]
\addplot[omThm,thick,mark=*,mark size=1pt] table[x=sd,y=cost]{figdata/strategies-2/25-the-central-risk-book/frontier.csv};
\addplot[gray,thick,dashed,mark=square*,mark size=1pt] table[x=sd_nofut,y=cost_nofut]{figdata/strategies-2/25-the-central-risk-book/frontier.csv};
\addplot[omDef,thick,dotted] coordinates {(429.9,0) (429.9,290)};
\legend{avec surcouche de futures sur indices,actions seules,limite de valeur en risque de 1~million de dollars}
\end{axis}
\end{tikzpicture}

\omcaption{Frontière coût-risque du central book selon la part de l'inventaire écoulée chaque jour, de 5~\% à 100~\%, avec et sans surcouche par des futures
sur indices~; la ligne pointillée représente l'écart type auquel la valeur en risque à un jour et à 99~\% atteint 1~million de dollars. Données~:
\texttt{s2\_crbook.frontier}.}\label{fig:s2:the-central-risk-book:frontier}
\end{omfigure}

\section{Gouvernance}

Un \omterm{def:s2:the-central-risk-book:crb}{central risk book} modifie qui détient le risque et qui paie la couverture. Les desks qui apportent les flux doivent se voir facturer un prix de transfert pour le risque
qu'ils remettent, et être crédités de la compensation que leurs flux procurent. Si le prix est trop bas, les desks transmettent un risque qu'ils auraient dû tarifer auprès de leurs clients~; si
il est trop élevé, ils couvrent de leur côté et la compensation est perdue. Les limites propres au book (valeur en risque, expositions aux facteurs, concentration, délai de
liquidation) requièrent une gestion des risques indépendante, et le croisement lui-même doit obéir à des règles. Les transactions des clients doivent être exécutées aux conditions promises, et
les informations issues des clients d'un desk ne doivent pas parvenir au desk de trading d'un autre. Le livre~7 (chapitre~27) traite du croisement interne des stratégies propres à une société~;
les mêmes questions d'allocation équitable se posent entre desks.

\section{Fiches de stratégie}

\begin{strategyfile}{Compensation centrale entre desks}\label{strat:s2:the-central-risk-book:netting}
\sfield{Qui paie, et pourquoi} La fourchette du marché et l'impact de marché, que ne rapportent pas les flux qui se compensent.
\sfield{Instruments et plateformes} Les positions des desks~; le moteur de croisement interne.
\sfield{Signal} Les expositions compensatoires entre desks.
\sfield{Dimensionnement et exécution} Calculer la position nette avant couverture, avec des prix de transfert vers les desks.
\sfield{Coûts} Systèmes et gouvernance.
\sfield{Comment elle meurt} Des flux corrélés qui ne se compensent pas~; des desks qui contournent le dispositif.
\sfield{Horizon, capacité, infrastructure} Intrajournalier~; un flux de positions à l'échelle de la société.
\sfield{Backtest honnête} Les flux réels des desks et leur corrélation, et non leur indépendance.
\sfield{Sources} Butz et Oomen (2019)~; ce chapitre~: 19.5~\% du coût de couverture économisé par la mutualisation de cinq desks.
\end{strategyfile}

\begin{strategyfile}{Risque résiduel couvert par facteurs}\label{strat:s2:the-central-risk-book:factor}
\sfield{Qui paie, et pourquoi} Le bon marché des contrats à terme sur indices par rapport aux actions prises isolément.
\sfield{Instruments et plateformes} Contrats à terme sur indices, paniers de facteurs, ETF sectoriels.
\sfield{Signal} Les expositions aux facteurs du book mutualisé, issues d'un modèle de risque.
\sfield{Dimensionnement et exécution} Couvrir les facteurs immédiatement~; laisser le risque spécifique.
\sfield{Coûts} Le roll et la base des contrats à terme.
\sfield{Comment elle meurt} Un modèle de risque qui omet un facteur.
\sfield{Horizon, capacité, infrastructure} Quotidien~; un modèle de risque factoriel (livre~7, chapitre~24).
\sfield{Backtest honnête} Les expositions aux facteurs hors échantillon.
\sfield{Sources} Ce chapitre~: un cinquième du risque porté éliminé pour 500~\$ par jour.
\end{strategyfile}

\begin{strategyfile}{Couverture patiente du risque mutualisé}\label{strat:s2:the-central-risk-book:patient}
\sfield{Qui paie, et pourquoi} L'impact évité en négociant de plus petites quantités et en laissant les flux ultérieurs se compenser.
\sfield{Instruments et plateformes} Les actions, négociées sur plusieurs jours.
\sfield{Signal} L'inventaire rapporté à la frontière coût-risque.
\sfield{Dimensionnement et exécution} Négocier chaque jour une part fixe de l'inventaire, fixée par la limite de risque.
\sfield{Coûts} Le risque résiduel.
\sfield{Comment elle meurt} Des flux à sens unique~; des pics de volatilité qui font dépasser la limite.
\sfield{Horizon, capacité, infrastructure} Des jours~; les limites de risque et le capital.
\sfield{Backtest honnête} Les flux et la volatilité des périodes de stress, et non des périodes moyennes.
\sfield{Sources} Garleanu et Pedersen (2013)~; ce chapitre~: 45\,300~\$ par jour pour une valeur en risque de 1~million de dollars, contre 335\,000~\$ desk par desk.
\end{strategyfile}

\section{Tutoriel~: couvrir ce qui reste}\label{tut:s2:the-central-risk-book:tut}

\begin{tutorial}
\textbf{Objectif.} Simuler les flux clients des desks, les couvrir desk par desk, de façon mutualisée et patiente avec une couche de contrats à terme, et choisir le taux de couverture à partir d'une
limite de risque. \textbf{État final~:} le tableau et deux figures.
\begin{tutsteps}
\item \textbf{Le book}.
\omcode{code/firm/crbook/firm_crbook.py}{62}{98}{Coût de couverture, desk par desk, mutualisé, et le central book patient.}\label{lst:s2:the-central-risk-book:book}
\item \textbf{La frontière et la limite}.
\omcode{code/strategies-2/25-the-central-risk-book/python/s2_crbook.py}{42}{59}{Coût et risque par taux~; le taux le moins coûteux dans une limite de valeur en risque.}\label{lst:s2:the-central-risk-book:frontier}
\item \textbf{Exécuter} \texttt{pooling()}, \texttt{frontier()}, \texttt{pick()}, \texttt{by\_desks()} et \texttt{fig\_crbook.py}.
\end{tutsteps}
\textbf{Pour aller plus loin.} Faire du flux d'un desk une couverture de celui d'un autre (corrélation négative)~; négocier les noms à des taux fixés par leur liquidité~;
ajouter un prix de transfert et voir quels desks contourneraient le book.
\end{tutorial}

\section{Construction~: le central risk book}\label{bld:s2:the-central-risk-book:crbook}

\begin{build}
\bfield{Rôle} Flux des desks, mutualisation, croisement interne, une couche de contrats à terme et couverture patiente, chiffrés avec le modèle du livre~7.
\bfield{Interface} \texttt{CRBConfig(\dots)}, \texttt{simulate\_desks(cfg)}, \texttt{hedge\_cost(trade, sim, cfg)}, \texttt{desk\_by\_desk},
\texttt{pooled}, \texttt{central(sim, cfg, rate, futures)}.
\bfield{Règles} Couvertures des desks exécutées comme si chacun était seul~; le central book négocie chaque jour une part fixe de l'inventaire~; des contrats à terme
couvrant l'inventaire pondéré par le bêta.
\bfield{Tests d'acceptation} \texttt{code/firm/crbook/tests/}~: des desks qui se compensent ne coûtent rien mutualisés~; négocier tout en même temps revient à mutualiser et
ne rapporte rien~; les contrats à terme éliminent le risque de marché~; le coût refait à la main.
\bfield{Pour aller plus loin} Taux nom par nom~; prix de transfert~; flux en période de stress.
\end{build}

\begin{omsources}
\item N.~Garleanu et L.~H.~Pedersen, ``Dynamic trading with predictable returns and transaction costs'', \emph{Journal of Finance} 68(6), 2013.
\item R.~Almgren, C.~Thum, E.~Hauptmann et H.~Li, ``Direct estimation of equity market impact'', \emph{Risk}, 2005.
\item M.~Butz et R.~Oomen, ``Internalisation by electronic FX spot dealers'', \emph{Quantitative Finance} 19(1), 2019.
\end{omsources}

\section{Exercices}

\begin{exercise}[$\star$]\label{exo:s2:the-central-risk-book:1}
Deux desks doivent acheter 3~millions de dollars et vendre 2~millions de dollars du même titre. Calculer le montant qui se croise en interne, et sa part dans le brut.
\end{exercise}

\begin{exercise}[$\star$]\label{exo:s2:the-central-risk-book:2}
Le P\&L porté a un écart type de 369\,000~\$ par jour. Calculer la valeur en risque à un jour à 99~\% sous hypothèse de normalité.
\end{exercise}

\begin{exercise}[$\star$]\label{exo:s2:the-central-risk-book:3}
Avec une demi-fourchette de 3~points de base, un impact $0.7\sigma\sqrt{q/V}$, $\sigma = 2~\%$ et $V = 50$ millions de dollars, combien coûte la couverture de 1~million de dollars~?
\end{exercise}

\begin{exercise}[$\star\star$]\label{exo:s2:the-central-risk-book:4}
Pourquoi la mutualisation fait-elle économiser moins que la part des flux qui se croisent~?
\end{exercise}

\begin{exercise}[$\star\star$]\label{exo:s2:the-central-risk-book:5}
Pourquoi ne négocier chaque jour qu'une partie de l'inventaire permet-il d'économiser autant de coût~?
\end{exercise}

\begin{exercise}[$\star\star$]\label{exo:s2:the-central-risk-book:6}
Pourquoi le gain de la mutualisation diminue-t-il au-delà de six desks lorsque les flux ont une partie commune~?
\end{exercise}

\begin{exercise}[$\star\star\star$]\label{exo:s2:the-central-risk-book:7}
\emph{Programmation.} Exécuter \texttt{pick(5e5)}. Quel taux une limite de valeur en risque de 500\,000~\$ retient-elle, et combien coûte-t-il~?
\end{exercise}

\begin{exercise}[$\star\star\star$]\label{exo:s2:the-central-risk-book:8}
\emph{Trouver la faille.} «~Le central book a réduit les coûts de couverture de 86~\%~; il faut négocier 5~\% par jour et les réduire de 94~\%.~»
\end{exercise}

\section{Problème~: couvrir ce qui reste}

\begin{problem}[{Problème du week-end --- le central risk book}]\label{pb:s2:the-central-risk-book:1}
Les desks synthétiques du chapitre et le relevé public.

\medskip\noindent\textbf{Partie I --- La mutualisation.}
\begin{enumerate}
\item Définir un \omterm{def:s2:the-central-risk-book:crb}{central risk book} et la \omterm{def:s2:the-central-risk-book:pooling}{mutualisation des risques}.
\item Décrire les desks, les flux, les rendements et les coûts.
\item Donner les coûts desk par desk et mutualisés, ainsi que la part croisée.
\item Quel choix Butz et Oomen ont-ils décrit~?
\end{enumerate}

\medskip\noindent\textbf{Partie II --- Le croisement.}
\begin{enumerate}[resume]
\item Comment l'économie dépend-elle du nombre de desks~?
\item Pourquoi la part commune compte-t-elle~?
\item Qu'ont établi Almgren et ses coauteurs au sujet de l'impact~?
\item Pourquoi l'économie est-elle plus faible que la part croisée~?
\end{enumerate}

\medskip\noindent\textbf{Partie III --- La couverture optimisée.}
\begin{enumerate}[resume]
\item Définir une \omterm{def:s2:the-central-risk-book:residual}{couverture du risque résiduel}.
\item Décrire la politique patiente et donner la frontière.
\item À quoi sert la superposition par futures~?
\item Quel taux une limite de 1~million de dollars retient-elle~?
\end{enumerate}

\medskip\noindent\textbf{Partie IV --- Le verdict.}
\begin{enumerate}[resume]
\item Énoncer le \emph{résultat nommé}~: le coût de couverture économisé par la mutualisation et le risque résiduel accepté.
\item Quelle hypothèse avantage la comparaison desk par desk~?
\item Que doit faire le prix de transfert alloué aux desks~?
\item Qu'est-ce qui peut mal tourner avec un livre patient dans un marché unidirectionnel~?
\item Que la gouvernance doit-elle tenir séparé~?
\item Quelle fiche de stratégie dépend d'un modèle de risque~?
\item Comment ce chapitre se rattache-t-il à l'internalisation du chapitre~24~?
\item En une phrase~: que vend un \omterm{def:s2:the-central-risk-book:crb}{central risk book} à la banque~?
\end{enumerate}
\end{problem}

\section{Questions d'entretien}

\begin{interviewq}[$\star$ \iqroles{trader}]\label{iq:s2:the-central-risk-book:1}
Qu'est-ce qu'un \omterm{def:s2:the-central-risk-book:crb}{central risk book}, et pourquoi une banque en gérerait-elle un~?
\end{interviewq}

\begin{interviewq}[$\star\star$ \iqroles{chercheur}]\label{iq:s2:the-central-risk-book:2}
Comment choisiriez-vous la vitesse à laquelle vous liquidez un inventaire mutualisé~?
\end{interviewq}

\begin{interviewq}[$\star\star$ \iqroles{trader}]\label{iq:s2:the-central-risk-book:3}
Deux desks vous envoient, sur la même action et au même moment, des risques opposés. Comment valorisez-vous le croisement vis-à-vis de chacun~?
\end{interviewq}

\begin{interviewq}[$\star\star$ \iqroles{risque}]\label{iq:s2:the-central-risk-book:4}
Quelles limites mettriez-vous sur un \omterm{def:s2:the-central-risk-book:crb}{central risk book}~?
\end{interviewq}

\begin{interviewq}[$\star\star$ \iqroles{développeur}]\label{iq:s2:the-central-risk-book:5}
Concevez le flux de positions à l'échelle de la firme dont un \omterm{def:s2:the-central-risk-book:crb}{central risk book} a besoin.
\end{interviewq}

\begin{interviewq}[$\star\star\star$ \iqroles{chercheur}]\label{iq:s2:the-central-risk-book:6}
$D$ desks ont des flux indépendants d'écart-type $s$ sur un même titre. Montrez que la mutualisation réduit le coût linéaire espéré d'un facteur $\sqrt D$ et
le coût d'impact espéré, proportionnel à $|q|^{3/2}$, d'un facteur $D^{1/4}$.
\end{interviewq}
